\section{Related Work}
\label{sec:related-work}

This section reviews prior work on scriptless GUI testing and TESTAR, including recent LLM-guided exploration, user stories as requirements artifacts and their structuring for automated analysis, and requirements traceability. It then positions the gap addressed by this research.

\subsection{Scriptless GUI Testing, TESTAR, and GUI State Models}
\label{subsec:scriptless-gui-models}

Scriptless GUI testing executes a system through its user interface without relying on manually written test scripts. Earlier work in GUI testing and model-based testing investigated event-flow graphs, widget hierarchies, and run-time interaction models to support automated test generation and coverage analysis~\citep{memon2007event}. TESTAR continues this line of research by dynamically exploring a System Under Test (SUT) through its GUI and recording observed behavior in a \emph{GUI State Model}~\citep{vos2021testar}. A GUI State Model represents the explored state space as a graph of states, actions, and transitions and provides detailed execution evidence for offline analysis. Prior work has used such evidence, for example, for offline oracles~\citep{de2019offline} and for detecting GUI changes across software versions~\citep{ricos2025delta}.

A key challenge is that the inferred state space can become large and highly interconnected. As exploration proceeds, state-model inference must deal with issues such as state abstraction, transition identification, and model growth. This challenge becomes more pronounced for larger systems, where model inference can be computationally demanding~\citep{mulders2022state,ricos2023distributed}. Consequently, GUI State Models are valuable technical artifacts, but they are also difficult to inspect and interpret manually.

A practical limitation of unguided exploration is that it does not necessarily exercise requirements-relevant workflows. To address this, recent work augments TESTAR with LLMs to support more goal-directed exploration, enabling interactions guided by high-level objectives such as user stories or Gherkin-like goals~\citep{van2025llm}. This is relevant when stakeholders care about exercising specific user workflows rather than only maximizing exploration breadth or reward. However, even when exploration is steered toward functional goals, the resulting GUI State Model is still expressed in low-level GUI terms, namely states, actions, and transitions. As a result, the problem of interpreting scriptless testing results in requirements terms remains unresolved.

\subsection{User Stories and Structured Scenario Representation}
\label{subsec:user-stories}

User stories are widely used in agile development to express functional requirements from a stakeholder perspective. A common template is the Connextra form popularized by Cohn: ``As a <role>, I want <goal> so that <reason>''~\citep{cohn2004user}. Although intentionally lightweight, user stories are typically written in natural language, which complicates automated analysis and tool support.

Prior work has sought to improve the quality and analyzability of user stories by imposing additional structure and quality criteria. The Quality User Story (QUS) framework, for example, defines criteria such as atomicity, unambiguity, and conceptual soundness to support the creation of clearer and more consistent stories~\citep{lucassen2016improving}. Related work has also shown that user stories can be transformed into more structured representations using natural language processing, for example by extracting actors, actions, and domain concepts for downstream analysis~\citep{lucassen2017visual,dalpiaz2018natural}. These approaches show that user stories can be made more analyzable without abandoning their role as lightweight requirements artifacts.

For this research, the key issue is not generating or assessing user stories in isolation, but representing them in a structured way that allows their essential steps or acceptance criteria to be recognized in execution evidence. This motivates the use of a structured representation that preserves requirements meaning while supporting comparison with GUI-level actions and paths.

\subsection{Requirements Traceability to Testing Artifacts}
\label{subsec:traceability}

Requirements traceability concerns establishing and maintaining links between requirements and other lifecycle artifacts such as design models, source code, and tests. The difficulty of sustaining such links as systems evolve has been recognized for decades; early work discusses the problem of ``requirements disconnects''~\citep{gotel1994analysis}. More recent surveys describe a broad range of traceability techniques, including information retrieval (IR), machine learning, and hybrid approaches that combine automated recovery with human validation~\citep{cleland2014software}.

A substantial portion of this literature focuses on links between requirements and testing artifacts. IR-based approaches, for example, attempt to recover links between textual requirements and manually authored test cases by exploiting lexical and semantic similarity~\citep{hayes2003improving,oliveto2010equivalence}. Such methods are useful when test artifacts contain explicit functional intent, domain vocabulary, or stable textual descriptions.

However, scriptless GUI testing produces a fundamentally different kind of artifact. Rather than manually specified test cases, it yields low-level execution evidence in the form of GUI states, actions, and transitions. These artifacts are tool-specific, graph-based, and often large, and they do not necessarily contain functional language that aligns directly with requirements. Consequently, existing requirement-to-test traceability techniques cannot be applied straightforwardly to TESTAR-generated GUI State Models.

\subsection{Gap and Positioning of This Research}
\label{subsec:positioning}

The reviewed literature shows strong progress in three relevant areas: generating execution evidence through scriptless GUI exploration, structuring user stories to improve analyzability, and recovering traceability links between requirements and traditional development and testing artifacts. However, these lines of work have not yet been integrated in a way that supports tracing requirements-level workflows to execution evidence captured in GUI State Models.

Even when TESTAR exploration is guided toward functional goals, the resulting evidence remains encoded as low-level GUI states and transitions. Stakeholders therefore face the same practical problem identified in the introduction: relevant behavior may be present in the model, but it is not expressed in requirements terms, and the scale of the model makes manual inspection impractical. This research addresses that gap by developing a method to represent user stories as a small set of essential steps or acceptance criteria, trace those elements to execution evidence in TESTAR-generated GUI State Models, and aggregate the resulting evidence into requirements coverage labels that can be reported to stakeholders.
